% Source: https://artofproblemsolving.com/community/c3h2478491_conjecture
% Stable identity: aops-topic-2478491-q01
\textbf{(a)} Since $a\ne b$, $\left|a^2-b^2\right|$ is a positive integer, and so are $2ab$ and $a^2+b^2$. Expanding,
\[\left(a^2-b^2\right)^2+(2ab)^2=a^4-2a^2b^2+b^4+4a^2b^2=a^4+2a^2b^2+b^4=\left(a^2+b^2\right)^2.\]

\textbf{(b)} With $a=2$, $b=1$ we get $3,4,5$, and with $a=1$, $b=4$ we get $15,8,17$: valid triples even though $b=1$ or $a=1$. But with $a=b$ (for example $a=b=2$) we get $\left|a^2-b^2\right|=0$, which is not the length of a side, so $a\ne b$ is necessary.

\textbf{(c)} For $z=a+bi$ we have $|z^2|=|z|^2$. Here $|z|^2=a^2+b^2$, while $z^2=\left(a^2-b^2\right)+2ab\,i$ has modulus $\sqrt{\left(a^2-b^2\right)^2+(2ab)^2}$. Equating gives the identity in (a): the real and imaginary parts of $z^2$ are the legs, and $|z|^2$ is the hypotenuse.
